\input{/home/pastor/.preamble_pl.tex}

\begin{document}
	
	\begin{center}
		\Large\textbf{Obsługa sieci sensorycznej/Zarządzanie siecią sensoryczną}
	\end{center}
	
	
	\tableofcontents
	
	
	
	\newpage
	\section{Ogólna idea oraz przeznaczenie układu}
	
	\section{Część elektroniczna}
		\subsection{Host (schemat blokowy)}
			\subsubsection{Moduł hosta 1}
			\subsubsection{Moduł hosta n}
			\subsubsection{Zasilanie}
			\subsubsection{Komunikacja}
			\subsubsection{...}
		
		\subsection{Płytki (\footnotesize{kilka płytek -- dla różnych przypadków) z czujnikami (schemat blokowy/schemat płytki w Eagle/KiCad)}}
			\subsubsection{Czujnik 1} 
			\subsubsection{Czujnik n}
			\subsubsection{Zasilanie (\footnotesize{można coś wymyślić z płytką zero-energetyczną -- energia pochodzi od drgań lub panelu fotowoltaicznego (DC/DC EHE004 lub BQ25504))}}
			\subsubsection{Komunikacja (bezprzewodowa \& przewodowa (I2C, SPI, ...))}
		
	
	\section{Część programistyczna}
		\subsection{Host (UML, schemat blokowy -- architektura publikacja-subskrypcja)}
			\subsubsection{Protokół}
			\subsubsection{Klasa 1}
			\subsubsection{Klasa n}
		
		\subsection{Płytka (UML, schemat blokowy -- współpraca z czujnikami/modułami)} 
			\subsubsection{Układy peryferyjne ESP32/Raspberry Pi z punktu widzenia softwearowego}
			\subsubsection{FiLtR KaLmAnA realizowany softwearowo}
			\subsubsection{Serializacja zebranych danych (nagłówek z opisem + dane)}
			
	\section{Montaż}
		\subsection{Wybrane elementy}
		\subsection{Osadzanie elementów (lutowanie/pasta/klejenie)}
		\subsection{Wydrukowane obudowy (układy zero energetyczne -- miejsce pod panel fotowoltaiczny w obudowie)} 
		
	
	\section{Pomiary/rozwiązane problemy/testy/...}
	
	\section{Przypadki użycia}
	
	\section{Wnioski} 
	
	
	\section{Dodatek -- zdjęcia/przebiegi/charakterystyki}
	
	
	
	
	
\end{document}